% Part: sets-functions-relations
% Chapter: relations
% Section: orders

\documentclass[../../../include/open-logic-section]{subfiles}

\begin{document}

\olfileid{sfr}{rel}{ord}
\olsection{క్రమాలు}

\begin{explain}
మనం చేసే అనేక పోలికలలో, ఒక నిర్దిష్ట దృష్టిలో కొన్ని వస్తువులు ఇతర
వస్తువుల ``కంటే చిన్నవి'', ``సమానమైనవి'', లేదా ``కంటే పెద్దవి''
అని చెబుతాం. వీటిలో \emph{క్రమ} సంబంధాలు ఇమిడి ఉంటాయి.
కానీ క్రమ సంబంధాలలో వేర్వేరు రకాలు ఉన్నాయి.
ఉదాహరణకు కొన్ని రకాలలో ఏ రెండు వస్తువులనైనా పోల్చగలగాలి; మరికొన్నింటిలో
అది అవసరం లేదు. కొన్ని తాదాత్మ్యాన్ని చేర్చుకుంటాయి ($\le$ వంటివి);
మరికొన్ని దానిని మినహాయిస్తాయి ($<$ వంటివి).
ఇక్కడ ఒక వర్గీకరణ మనకు ఉపయోగపడుతుంది.
\end{explain}

\begin{defn}[పూర్వక్రమం]
స్వావర్తన, సంక్రామక ధర్మాలు రెండూ ఉన్న సంబంధాన్ని
\emph{పూర్వక్రమం} అంటాం.
\end{defn}

\begin{defn}[పాక్షిక క్రమం]
ప్రతిసౌష్ఠవ ధర్మం కూడా ఉన్న పూర్వక్రమాన్ని
\emph{పాక్షిక క్రమం} అంటాం.
\end{defn}

\begin{defn}[రేఖీయ క్రమం]\ollabel{def:linearorder}
సంయుక్త ధర్మం కూడా ఉన్న పాక్షిక క్రమాన్ని
\emph{సంపూర్ణ క్రమం} లేదా \emph{రేఖీయ క్రమం} అంటాం.
\end{defn}

\begin{ex}
ప్రతి రేఖీయ క్రమమూ పాక్షిక క్రమమే; ప్రతి పాక్షిక క్రమమూ పూర్వక్రమమే.
కానీ తిరుగు వాదనలు చెల్లవు. $A$పై సార్వత్రిక సంబంధం స్వావర్తన,
సంక్రామక ధర్మాలు కలిగి ఉంటుంది కనుక అది పూర్వక్రమం.
అయితే $A$లో ఒకటి కంటే ఎక్కువ !!{element} ఉంటే, సార్వత్రిక
సంబంధానికి ప్రతిసౌష్ఠవ ధర్మం ఉండదు; కనుక అది పాక్షిక క్రమం కాదు.
\end{ex}

\begin{ex}
$\Bin^*$పై \emph{కంటే పొడవైనది కాదు} అనే $\preccurlyeq$ సంబంధాన్ని
పరిశీలించండి: $x \preccurlyeq y$ కావడానికి అవసరమైన మరియు సరిపడే
షరతు $\len{x} \le \len{y}$ అవడం.
ఇది పూర్వక్రమం (స్వావర్తన, సంక్రామక ధర్మాలు ఉన్నాయి);
సంయుక్త ధర్మం కూడా ఉంది. కానీ ప్రతిసౌష్ఠవ ధర్మం లేదు కనుక
ఇది పాక్షిక క్రమం కాదు. ఉదాహరణకు $01 \preccurlyeq 10$,
$10 \preccurlyeq 01$ అయినా $01 \neq 10$.
\end{ex}

\begin{ex}
సమితుల సమితిపై $\subseteq$ సంబంధం ఒక ముఖ్యమైన పాక్షిక క్రమం.
సాధారణంగా ఇది రేఖీయ క్రమం కాదు. ఎందుకంటే $a \neq b$ అనుకొని
$\Pow{\{a, b\}} = \{\emptyset, \{a\}, \{b\}, \{a,b\}\}$ను పరిశీలిస్తే,
$\{a\} \nsubseteq \{b\}$, $\{a\} \neq \{b\}$,
$\{b\} \nsubseteq \{a\}$ అని తెలుస్తుంది.
\end{ex}

\begin{ex}
\emph{శేషం లేకుండా భాగించడం} అనే సంబంధం, రేఖీయ క్రమం కాని
పాక్షిక క్రమానికి ఉదాహరణ. పూర్ణ సంఖ్యలు $n$, $m$లకు,
$n$ అనేది $m$ను శేషం లేకుండా భాగిస్తుందని చెప్పడానికి $n \mid m$
అని రాస్తాం. అంటే ఏదో పూర్ణ సంఖ్య $k$కు $m = kn$ అయితే, అయితేనే.
$\Nat$పై ఇది పాక్షిక క్రమం, కానీ రేఖీయ క్రమం కాదు:
ఉదాహరణకు $2 \nmid 3$, అలాగే $3 \nmid 2$.
$\Int$పై సంబంధంగా పరిగణిస్తే, భాగించడం కేవలం పూర్వక్రమమే;
దానికి ప్రతిసౌష్ఠవ ధర్మం లేదు. ఎందుకంటే $1 \mid -1$,
$-1 \mid 1$ అయినా $1 \neq -1$.
\end{ex}

\begin{ex}
క్రమాల సమితి $A^*$పై \emph{పొడిగింపు} సంబంధం ఇలా ఉంటుంది:
$s \sqsubseteq s'$ కావడానికి అవసరమైన మరియు సరిపడే షరతు
$s = \emptyseq$ (శూన్య క్రమం) కావడం, $s = s'$ కావడం, లేదా
$s = \tuple{s_1, \dots, s_n}$,
$s' = \tuple{s_1, \dots, s_n, s_{n+1}, \dots, s_m}$ కావడం.
$s \sqsubseteq s'$ అయితే, $s$ను $s'$ యొక్క \emph{ప్రారంభ ఖండం}
అని కూడా అంటాం. $A^*$పై పొడిగింపు సంబంధం పాక్షిక క్రమం,
కానీ రేఖీయ క్రమం కాదు. ఉదాహరణకు $a \neq b$ అయితే
$ab \not\sqsubseteq ba$, $ba \not\sqsubseteq ab$.
\end{ex}

\begin{defn}[కఠిన క్రమం]
అస్వావర్తన, ఏకదిశ, సంక్రామక ధర్మాలు ఉన్న సంబంధాన్ని
\emph{కఠిన క్రమం} అంటాం.
\end{defn}

\begin{defn}[కఠిన రేఖీయ క్రమం]\ollabel{def:strictlinearorder}
సంయుక్త ధర్మం కూడా ఉన్న కఠిన క్రమాన్ని
\emph{కఠిన సంపూర్ణ క్రమం} లేదా \emph{కఠిన రేఖీయ క్రమం} అంటాం.
\end{defn}

\begin{ex}
కఠిన రేఖీయ క్రమమైన $<$కు అనుగుణమైన రేఖీయ క్రమం $\le$.
కఠిన క్రమమైన $\subsetneq$కు అనుగుణమైన పాక్షిక క్రమం $\subseteq$.
\end{ex}

$A$పై ఏ కఠిన క్రమం $R$నైనా, వికర్ణమైన $\Id{A}$ను,
అంటే $\tuple{x, x}$ జతలన్నింటినీ చేర్చి పాక్షిక క్రమంగా మార్చవచ్చు.
(దీనిని $R$ యొక్క \emph{స్వావర్తన సంవృతం} అంటాం.)
తిరుగుగా, పాక్షిక క్రమం నుంచి $\Id{A}$ను తొలగించి కఠిన క్రమాన్ని
పొందవచ్చు. తదుపరి రెండు ఫలితాలు దీనిని కచ్చితంగా చెబుతాయి.

\begin{prop}\ollabel{prop:stricttopartial}
$R$ అనేది $A$పై కఠిన క్రమం అయితే, $R^+ = R \cup \Id{A}$
పాక్షిక క్రమం. అంతేకాక, $R$ కఠిన రేఖీయ క్రమం అయితే,
$R^+$ రేఖీయ క్రమం.
\end{prop}

\begin{proof}
$R$ కఠిన క్రమం అనుకుందాం. అంటే $R \subseteq A^2$, అలాగే
$R$కు అస్వావర్తన, ఏకదిశ, సంక్రామక ధర్మాలు ఉన్నాయి.
$R^+ = R \cup \Id{A}$ తీసుకుందాం.
$R^+$కు స్వావర్తన, ప్రతిసౌష్ఠవ, సంక్రామక ధర్మాలు ఉన్నాయని చూపాలి.

ప్రతి $x \in A$కూ $\tuple{x, x} \in \Id{A} \subseteq R^+$
కనుక, $R^+$కు స్వావర్తన ధర్మం స్పష్టంగా ఉంది.

$R^+$కు ప్రతిసౌష్ఠవ ధర్మం ఉందని చూపడానికి, విరోధ నిరూపణ కోసం
$R^+xy$, $R^+yx$ అయినా $x \neq y$ అనుకుందాం.
$\tuple{x,y} \in R \cup \Id{A}$ అయినా
$\tuple{x, y} \notin \Id{A}$ కనుక,
$\tuple{x, y} \in R$, అంటే $Rxy$, కావాలి.
ఇదే విధంగా $Ryx$. కానీ ఇది $R$ ఏకదిశ సంబంధమనే ఊహకు విరుద్ధం.

సంక్రామకత్వాన్ని స్థాపించడానికి $R^+xy$, $R^+yz$ అనుకుందాం.
$\tuple{x, y} \in R$, $\tuple{y,z} \in R$ రెండూ ఉంటే,
$R$ సంక్రామక సంబంధం కనుక $\tuple{x, z} \in R$.
లేకపోతే, $\tuple{x, y} \in \Id{A}$, అంటే $x = y$;
లేదా $\tuple{y, z} \in \Id{A}$, అంటే $y = z$.
మొదటి సందర్భంలో ఊహ ప్రకారం $R^+yz$ ఉంది, $x = y$ కనుక $R^+xz$.
రెండవ సందర్భమూ ఇలాగే. ఏ సందర్భంలోనైనా $R^+xz$;
అందువల్ల $R^+$కు సంక్రామక ధర్మం కూడా ఉంది.

``అంతేకాక'' అనే భాగం కోసం, $R$కు సంయుక్త ధర్మం కూడా ఉందనుకుందాం.
కనుక ప్రతి $x \neq y$కూ $Rxy$ లేదా $Ryx$;
అంటే $\tuple{x, y} \in R$ లేదా $\tuple{y, x} \in R$.
$R \subseteq R^+$ కనుక ఇది $R^+$కూ నిజమే.
అందువల్ల $R^+$కు సంయుక్త ధర్మం కూడా ఉంది.
\end{proof}

\begin{prop}\ollabel{prop:partialtostrict}
$R$ అనేది $A$పై పాక్షిక క్రమం అయితే,
$R^- = R \setminus \Id{A}$ కఠిన క్రమం.
అంతేకాక, $R$ రేఖీయ క్రమం అయితే $R^-$ కఠిన రేఖీయ క్రమం.
\end{prop}

\begin{proof}
దీనిని అభ్యాసంగా వదిలివేస్తున్నాం.
\end{proof}

\begin{prob}
\olref[sfr][rel][ord]{prop:partialtostrict}కు నిరూపణ ఇవ్వండి.
\end{prob}

కఠిన రేఖీయ క్రమాలకు మూలకాధారిత సమానత్వాన్ని పోలిన ధర్మం ఉంటుందని
ఈ సరళమైన ఫలితం స్థాపిస్తుంది:

\begin{prop}\ollabel{prop:extensionality-strictlinearorders}
$<$ అనేది $A$పై కఠిన రేఖీయ క్రమం అయితే:
\[
  (\forall a, b \in A)((\forall x \in A)(x < a \liff x < b) \lif a = b).
\]
\end{prop}

\begin{proof}
$(\forall x \in A)(x < a \liff x < b)$ అనుకుందాం.
$a < b$ అయితే $a < a$; ఇది $<$కు అస్వావర్తన ధర్మం ఉండటానికి
విరుద్ధం. కనుక $a \nless b$.
ఇదే విధంగా $b \nless a$. $<$కు సంయుక్త ధర్మం ఉంది కనుక $a = b$.
\end{proof}

\end{document}
